% !TeX root = ../../Book2.tex
% 纲要标题：### 20.2 Weyl 代数
% 纲要位置：计划纲要.md，第 1704 行。

\section{Weyl 代数}
\label{b2:sec:2002}

在多项式上，先乘不定元再求导，与先求导再乘不定元，恰好相差恒等算子。Weyl 代数只保留这条交换关系。我们先证明抽象代数中的正规单项式确实独立，再判断它对通常多项式空间的作用是否忠实；这两项问题在正特征中有不同答案。

% 纲要要点（供目录核对）：
%
% * 微分与乘法的交换关系；区分任意特征下的抽象 Weyl 代数、通常多项式作用与辅助变量忠实作用
% * 正规单项式及作用；下降重排给生成性，辅助算子在一上的取值记录两个指数而给独立性
% * 特征零与正特征的差别；中心秩 p²、有限维表示迹障碍和中心零纤维的矩阵商在正文完整证明；原题改为商模分解、平移纤维及中心分式除代数
%

\subsection{微分与乘法的交换关系}
\label{b2:subsec:2002:01}

在 \(k[X]\) 上令 \(X\) 表示乘以不定元的算子，令 \(D_0\) 为形式求导。Leibniz 公式\footnote{莱布尼茨（Gottfried Wilhelm Leibniz，1646-1716），德国数学家及哲学家，创立微积分，并研究符号演算。}给出 \(D_0X-XD_0=1\)，这在任意特征下都成立。

\begin{definition}\label{b2:def:first-weyl-algebra}
设 \(k\) 为域。将自由结合代数 \(k\langle x,\mathrm{d}\rangle\) 对关系 \(\mathrm{d}\cdot x-x\cdot\mathrm{d}=1\) 生成的双边理想取商，所得代数称为 \(k\) 上的第一\term[Weyl-daishu]{Weyl 代数}[first Weyl algebra]，记为
\[
A_1(k)=k\langle x,\mathrm{d}\rangle/(\mathrm{d}\cdot x-x\cdot\mathrm{d}-1).
\]
其代数同态由生成元的像决定，并要求这些像满足 \(\mathrm{d}\cdot x-x\cdot\mathrm{d}=1\)。
\end{definition}
\symidx{\(A_1(k)\)：第一 Weyl 代数}

由\cref{b2:prop:algebra-relations-universal}，\(x\mapsto X,\mathrm{d}\mapsto D_0\) 给出 \(A_1(k)\) 在 \(k[X]\) 上的作用。但关系成立只证明作用存在，尚未证明这个作用单射。

\subsection{正规单项式及其作用}
\label{b2:subsec:2002:02}

\begin{theorem}\label{b2:thm:weyl-normal-basis}
设 \(k\) 为域，\(A_1(k)=k\langle x,\mathrm{d}\rangle/(\mathrm{d}\cdot x-x\cdot\mathrm{d}-1)\)，则全部 \(x^i\mathrm{d}^j\)、\(i,j\ge0\)，构成 \(A_1(k)\) 的一组 \(k\) 基。因此每个元素唯一写为有限和 \(\sum_j f_j(x)\mathrm{d}^j\)，其中 \(f_j\in k[x]\)。
\end{theorem}
\begin{proof}
生成性来自 \(\mathrm{d}\cdot x=x\cdot\mathrm{d}+1\)。对一个字，按字长及其中位于某个 \(x\) 左边的 \(\mathrm{d}\) 的配对数作字典式归纳：将相邻 \(\mathrm{d}\cdot x\) 换成 \(x\cdot\mathrm{d}+1\)，第一项的字长不变、错序数减少一，第二项的字长减少二。这两个量都是非负整数，故不能无限下降，最终只剩全部 \(x\) 在前、全部 \(\mathrm{d}\) 在后的正规字。不过归约终止这里只保证正规字生成，尚不能排除它们之间仍有线性关系。

为证独立性，我们要找到一个作用，使 \(x^i\mathrm{d}^j\) 的像保留两个指数的信息。通常的求导算子在常数一上反复作用都为零，不能用一来检测 \(j\)。因此在两个交换变量的多项式空间 \(k[X,Y]\) 上令 \(X\) 仍表示乘法，并令
\[
D=\frac{\partial}{\partial X}+Y,
\]
其中 \(Y\) 也表示乘法算子。增加的乘 \(Y\) 算子与乘 \(X\) 交换，故仍有 \(DX-XD=1\)，给出一个表示。同时 \(\partial(Y^j)/\partial X=0\)，所以 \(D(Y^j)=Y^{j+1}\)，从常数一出发便有 \(D^j(1)=Y^j\)，因而
\[
X^iD^j(1)=X^iY^j.
\]
若抽象代数中有有限关系 \(\sum_{i,j}c_{ij}x^i\mathrm{d}^j=0\)，令该算子作用在一上，得到 \(\sum_{i,j}c_{ij}X^iY^j=0\)。交换多项式的单项式独立，故所有 \(c_{ij}=0\)。整个论证没有使用阶乘可逆性。
\end{proof}

这个双变量表示因而在任意特征下都忠实。在特征 \(p\) 中，虽然通常的第 \(p\) 次求导在 \(k[X]\) 上为零，这里的 \(D^p(1)=Y^p\) 却不为零。它证明的是抽象 Weyl 代数中正规字的独立性，并不能据此断言通常的多项式作用也忠实；后者将在\cref{b2:prop:weyl-polynomial-action}中单独判断。后面的\cref{b2:thm:diamond-lemma}还会把改写的下降条件与歧义检查结合起来，证明正规式唯一。Weyl 代数的正规单项式基还可参阅 Etingof 等人\cite[Proposition 2.7.1(i), pp. 17--18]{EtingofEtAl2011RepresentationTheory}；那里使用另一种带辅助变量的表示。

\ProseParagraphBreak
用字长滤过，\(F_nA_1(k)\) 恰由 \(i+j\le n\) 的正规单项式生成。交换关系中的修正项降低次数二，故最高次部分变得交换。

\begin{proposition}\label{b2:prop:weyl-associated-graded}
设 \(k\) 为域，\(A_1(k)=k\langle x,\mathrm{d}\rangle/(\mathrm{d}\cdot x-x\cdot\mathrm{d}-1)\)，按 \(x,\mathrm{d}\) 的字长赋予滤过 \(F_nA_1(k)\)，则 \(\operatorname{gr}_F A_1(k)\cong k[X,D]\)。因此 \(A_1(k)\) 没有零因子，且左右 Noether。
\end{proposition}
\begin{proof}
将交换变量 \(X,D\) 分别送到一次主符号 \(\operatorname{symb}(x),\operatorname{symb}(\mathrm{d})\)。它们交换且生成伴随分次，故得到满同态。\cref{b2:thm:weyl-normal-basis}说明第 \(n\) 个商片恰以 \(i+j=n\) 的单项式符号为基，因此同态逐次单射。

交换多项式环 \(k[X,D]\) 是整环；由第一卷定理12.4.10（Hilbert 基定理），它还为 Noether 环。应用\cref{b2:prop:filtered-domain}及\cref{b2:thm:filtered-noetherian-transfer}，得到无零因子性和两侧 Noether 性。
\end{proof}

例如 \(\mathrm{d}^2x^3=x^3\mathrm{d}^2+6x^2\mathrm{d}+6x\)，由两次使用 \(\mathrm{d}\cdot x^n=x^n\mathrm{d}+nx^{n-1}\) 得到。等式在正特征中仍成立，只是整数系数按该特征取值。

\subsection{特征零与正特征}
\label{b2:subsec:2002:03}

为避免与群交换子混淆，本节的 \([a,b]=ab-ba\) 表示代数中的加法交换子。直接归纳给出
\begin{equation}\label{b2:eq:weyl-commutators}
[\mathrm{d},x^i\mathrm{d}^j]=i x^{i-1}\mathrm{d}^j,\qquad
[x,x^i\mathrm{d}^j]=-j x^i\mathrm{d}^{j-1}.
\end{equation}
零指数时右侧解释为零。

\begin{theorem}\label{b2:thm:weyl-center-simplicity}
设 \(k\) 为域，\(A_1(k)=k\langle x,\mathrm{d}\rangle/(\mathrm{d}\cdot x-x\cdot\mathrm{d}-1)\)。若 \(\operatorname{char}k=0\)，则 \(Z\bigl(A_1(k)\bigr)=k\)，且 \(A_1(k)\) 是单环。若 \(\operatorname{char}k=p>0\)，则
\[
Z\bigl(A_1(k)\bigr)=k[x^p,\mathrm{d}^p],
\]
这是二元多项式环；\(A_1(k)\) 是其中心上的自由模，一组基为
\[
\{x^i\mathrm{d}^j:0\le i,j<p\},
\]
故在中心上秩为 \(p^2\)。此时 \(A_1(k)\) 不是单环。
\end{theorem}
\begin{proof}
由\eqref{b2:eq:weyl-commutators}及正规单项式基，元素 \(\sum c_{ij}x^i\mathrm{d}^j\) 同时与 \(x,\mathrm{d}\) 交换，当且仅当每个非零系数满足 \(i c_{ij}=j c_{ij}=0\)。在特征零中只剩 \(i=j=0\)；在特征 \(p\) 中要求 \(p\mid i,j\)。这证明两个中心公式，且正规单项式基保证 \(x^p,\mathrm{d}^p\) 之间没有多项式关系。

在特征零中，要证明非零双边理想 \(I\) 等于整个代数，只需在其中得到非零标量。双边性保证与 \(x\) 或 \(\mathrm{d}\) 取交换子后仍留在 \(I\)，而\eqref{b2:eq:weyl-commutators}表明这分别降低 \(\mathrm{d}\) 或 \(x\) 的指数。选其中非零元素 \(P=\sum_{j=0}^m f_j(x)\mathrm{d}^j\)，使 \(\mathrm{d}\) 次数 \(m\) 最小。若 \(m>0\)，则 \([P,x]\in I\) 的 \(\mathrm{d}\) 次数为 \(m-1\)，最高系数为 \(m\cdot f_m(x)\ne0\)，矛盾。因此 \(I\) 含非零 \(f(x)=a_rx^r+\cdots+a_0\)，其中 \(a_r\ne0\)。记 \(\operatorname{ad}_{\mathrm{d}}(u)=[\mathrm{d},u]\)，则
\[
\operatorname{ad}_{\mathrm{d}}^{\,r}(f)=r!a_r\in k^\times.
\]
这个交换子仍属于 \(I\)，故 \(1\in I\)，\(I=A_1(k)\)。这里 \(r!\ne0\) 正是特征零假设的作用。

在特征 \(p\) 中，将每个指数唯一写为 \(i=pa+r\)、\(j=pb+s\)，其中 \(0\le r,s<p\)。利用 \(x^p,\mathrm{d}^p\) 居中，便有
\[
x^i\mathrm{d}^j=(x^p)^a(\mathrm{d}^p)^b x^r\mathrm{d}^s,
\]
所以这些余数指数的单项式在中心上生成 \(A_1(k)\)。若它们之间有中心系数的线性关系，将各系数展开为 \(x^p,\mathrm{d}^p\) 的多项式，不同商和余数给出不同的正规单项式。\cref{b2:thm:weyl-normal-basis}迫使所有系数为零，证明中心上的自由性及秩公式。中心元 \(x^p\) 生成的双边理想就是 \(x^pA_1(k)\)，按正规单项式基它由 \(x\) 指数至少为 \(p\) 的单项式生成。因此它非零且不含一，是一个真双边理想。
\end{proof}

\begin{proposition}\label{b2:prop:weyl-polynomial-action}
设 \(k\) 为域，\(A_1(k)=k\langle x,\mathrm{d}\rangle/(\mathrm{d}\cdot x-x\cdot\mathrm{d}-1)\)。令 \(x\) 在 \(k[X]\) 上作用为乘 \(X\)，\(\mathrm{d}\) 作用为形式求导，则这个 \(A_1(k)\) 模作用在 \(\operatorname{char}k=0\) 时忠实，在 \(\operatorname{char}k>0\) 时不忠实。
\end{proposition}
\begin{proof}
特征零时，若非零 \(P=\sum_{j=0}^m f_j(x)\mathrm{d}^j\) 作用为零，记 \(M_X\) 为乘 \(X\) 的算子。由\eqref{b2:eq:weyl-commutators}，\([f_j(X)\mathrm{d}^j,M_X]=j f_j(X)\mathrm{d}^{j-1}\)。因此从 \(P\) 开始反复按 \([P,M_X]\) 的顺序取 \(m\) 次交换子，仍作用为零；低于 \(m\) 次的项全部消失，所得算子是乘 \(m!f_m(X)\)，令其作用在一上便得到矛盾。\(m=0\) 时直接如此。

特征 \(p\) 时，第 \(p\) 次形式求导在每个 \(X^n\) 上都为零，因为连续 \(p\) 个整数的乘积含因子 \(p\)。所以非零的抽象元素 \(\mathrm{d}^p\) 作用为零，非零性则由\cref{b2:thm:weyl-normal-basis}保证。
\end{proof}

\begin{proposition}\label{b2:prop:weyl-finite-dimensional-representations}
设 \(k\) 为任意域，\(A=A_1(k)=k\langle x,\mathrm{d}\rangle/(\mathrm{d}\cdot x-x\cdot\mathrm{d}-1)\)。若 \(\operatorname{char}k=0\)，则每个有限维幺性左 \(A\) 模都是零模。若 \(\operatorname{char}k=p>0\)，则每个有限维幺性左 \(A\) 模的 \(k\) 维数都被 \(p\) 整除；此时
\[
V_p=k[X]/(X^p)
\]
在 \(x\) 作用为乘 \(X\)、\(\mathrm{d}\) 作用为形式求导时，是一个 \(p\) 维简单左 \(A\) 模，该作用诱导 \(k\) 代数同构
\[
A/(x^p,\mathrm{d}^p)\cong\operatorname{End}_k(V_p)\cong M_p(k),
\]
其中 \((x^p,\mathrm{d}^p)\) 是这两个中心元生成的双边理想。
\end{proposition}
\begin{proof}
在一个有限维幺性左 \(A\) 模 \(V\) 上，记 \(x,\mathrm{d}\) 的作用为 \(X,D\)，则 \(DX-XD=1_V\)。矩阵迹满足 \(\operatorname{Tr}(DX)=\operatorname{Tr}(XD)\)，故
\[
0=\operatorname{Tr}(DX-XD)=(\dim_kV)\,1_k.
\]
在特征零中这迫使 \(\dim_kV=0\)；在特征 \(p\) 中则迫使 \(p\mid\dim_kV\)。这一迹障碍也见 Etingof 等人\cite[Problem 2.7.4, p. 18]{EtingofEtAl2011RepresentationTheory}。

以下设 \(\operatorname{char}k=p>0\)。形式求导保持理想 \((X^p)\)，因为 \((X^pf)'=X^pf'\)，所以乘法与求导都下降到 \(V_p\)，仍满足 Weyl 关系。若一个非零子模含有非零 \(f\)，取次数小于 \(p\) 的代表，设其次数为 \(r\)、最高系数为 \(a_r\ne0\)。反复求导 \(r\) 次得到非零常数 \(r!a_r\)；这里 \(r<p\)，所以阶乘不为零。因此子模含一，再乘 \(X\) 就得到 \(1,X,\ldots,X^{p-1}\)，故 \(V_p\) 简单。

在 \(V_p\) 上，\(X^p\) 与 \(D^p\) 都为零，故作用经过所写商代数。要证明其像是全部线性算子，可以在像中逐一构造矩阵单位。令 \(H=XD\)，则 \(H(X^r)=rX^r\) 对 \(0\le r<p\)。因而
\[
E_0=\prod_{r=1}^{p-1}\frac{H-r\,1_{V_p}}{-r}
\]
恰为到常数项的投影：它在一上为恒等，在其余基向量上为零。所有分母都是 \(k\) 中的非零元素。对 \(0\le i,j<p\)，令
\[
E_{ij}=\frac{1}{j!}X^iE_0D^j.
\]
当 \(r<j\) 时，\(D^j(X^r)=0\)；当 \(r=j\) 时，它等于 \(j!\)；当 \(r>j\) 时，它是正次数单项式的标量倍，随后被 \(E_0\) 消去。因此 \(E_{ij}(X^r)=\delta_{jr}X^i\)，正是该基下的矩阵单位。它们全在商代数的像中，故诱导映射满射。另一方面，商代数由 \(x^i\mathrm{d}^j\)、\(0\le i,j<p\) 张成，维数至多 \(p^2\)；目标维数恰为 \(p^2\)，所以这个满射也是单射。
\end{proof}
